\chapter{Среда работы и воспроизводимый контур}

\section{Назначение главы}

Первая глава вводит читателя в рабочую среду курса и объясняет, почему анализ данных следует рассматривать как воспроизводимый вычислительный процесс, а не как набор изолированных скриптов и файлов.

\section{Jupyter как рабочая среда инженерного исследования}

В рамках настоящего пособия Jupyter-блокнот понимается как вычислительный документ, объединяющий код, текстовые пояснения, таблицы, графики и результаты вычислений.
Такой формат особенно удобен для учебного процесса и инженерного анализа по следующим причинам:
\begin{enumerate}[leftmargin=1.25cm]
\item ход вычислений остается открытым и может быть проверен по шагам;
\item теоретическое объяснение и программная реализация находятся в одном артефакте;
\item промежуточные результаты могут сопровождаться интерпретацией непосредственно после вычисления;
\item материал легко использовать как для демонстрации преподавателем, так и для самостоятельной работы обучающегося.
\end{enumerate}

В качестве основной среды в пособии рассматривается JupyterLab.
При необходимости те же файлы формата \code{.ipynb} могут быть открыты и в Jupyter Notebook.
С методической точки зрения это различие вторично: существенным является не конкретная оболочка, а воспроизводимость вычислений и сохранение связки ``код -- комментарий -- результат''.

\screenshotplaceholder
  {fig:jupyterlab-start}
  {Главное окно JupyterLab}
  {Стартовое окно JupyterLab с файловой панелью, списком каталогов проекта и открытым учебным блокнотом.}
  {Открыть JupyterLab из установленного учебного ядра. Сделать скриншот главного окна так, чтобы были видны файловая панель, верхнее меню и имя открытого блокнота \code{01\_course\_setup.ipynb}.}

\figureinstruction
  {Главное окно JupyterLab, используемое как основная рабочая среда учебного комплекта.}
  {После рисунка необходимо пояснить, что для дальнейшей работы студенту важно видеть не только сам блокнот, но и структуру каталогов, поскольку данные, ноутбуки, главы пособия и вспомогательный код разнесены по разным каталогам.}

\section{Структура учебного комплекта}

Воспроизводимость обеспечивается не одним блокнотом, а всей структурой учебного комплекта.
В репозитории выделяются:
\begin{itemize}[leftmargin=1.25cm]
\item \code{docs/} --- текстовые и методические материалы;
\item \code{notebooks/} --- исполняемые учебные сценарии;
\item \code{data/raw/}, \code{data/sample/}, \code{data/processed/} --- уровни работы с данными;
\item \code{src/} --- повторно используемые функции;
\item \code{kernel/} --- описание учебного ядра и установочные скрипты.
\end{itemize}

Такое разделение необходимо для того, чтобы:
\begin{enumerate}[leftmargin=1.25cm]
\item сохранять исходные данные отдельно от подготовленных;
\item выносить громоздкие функции из блокнотов;
\item обеспечивать единый состав библиотек;
\item связывать главы пособия с соответствующими блокнотами и данными.
\end{enumerate}

\screenshotplaceholder
  {fig:repo-structure}
  {Структура каталогов учебного репозитория}
  {Окно проводника файлов или JupyterLab с раскрытой структурой каталогов \code{docs}, \code{notebooks}, \code{data}, \code{src}, \code{kernel}.}
  {Сделать скриншот структуры репозитория так, чтобы были видны названия основных каталогов. Желательно использовать единый масштаб и выделить корень проекта.}

\figureinstruction
  {Структура учебного репозитория и распределение материалов по каталогам.}
  {После рисунка следует пояснить, что данная структура не является формальной избыточностью: она необходима для разделения теоретического текста, вычислительных сценариев, данных и повторно используемого кода.}

\section{Учебное ядро}

Для исключения расхождений в версиях библиотек и способе запуска блокнотов в комплект включено отдельное учебное ядро Python.
Установочный сценарий создает виртуальное окружение, устанавливает необходимые зависимости и регистрирует kernel \code{Power Data Analysis} в системе Jupyter.

Это решение важно методически и практически.
Во-первых, оно снижает порог входа для обучающегося, которому не требуется вручную собирать окружение.
Во-вторых, оно делает воспроизводимыми примеры, графики и вычисления, приведенные в пособии.

\section{Что должен сделать обучающийся}

После изучения данной главы и выполнения связанного блокнота обучающийся должен:
\begin{itemize}[leftmargin=1.25cm]
\item установить учебное ядро;
\item открыть JupyterLab или Jupyter Notebook;
\item запустить установочный блокнот;
\item убедиться в доступности основных библиотек;
\item проверить, что структура каталогов читается корректно.
\end{itemize}
